\section{Asian and scope of the instrument}

This appendix reproduces the research instruments of the AI-Driven Vulnerability
Analysis Engine (AVAE, Component~4): the prepared data set, the feature set, the
model configuration, the evaluation metrics, and the component artefacts.

\section{Prepared data set}

The primary data set is a curated collection of 155,852 vulnerability records
derived from public vulnerability sources. Class~0 (no recorded CISA KEV
membership) contains 154,889 records and class~1 (recorded KEV membership)
contains 963 records, which is a heavily imbalanced distribution. The class
labels identify recorded KEV membership and are used as the binary target for
the classifier. Supporting intelligence is drawn from the National
Vulnerability Database (NVD), the CISA Known Exploited Vulnerabilities (KEV)
catalogue, and the Exploit Prediction Scoring System (EPSS).

\section{Feature engineering}

The classifier uses five numerical features:

\begin{enumerate}
  \item CVSS v3 score -- severity of the vulnerability;
  \item EPSS score -- probability of exploitation as estimated by EPSS;
  \item EPSS percentile -- rank of the EPSS score within the distribution;
  \item vulnerability age -- time since the vulnerability was published;
  \item CVSS--EPSS interaction -- the product of the CVSS v3 score and the EPSS
  score, capturing the joint effect of severity and exploitation likelihood.
\end{enumerate}

Data cleaning removed duplicates, handled invalid values, converted attributes
to numerical formats, and computed vulnerability age from publication dates.

\section{Model configuration and evaluation}

XGBoost was selected for its support of nonlinear relationships and efficient
learning from structured numerical data. Candidate models were compared through
hyperparameter tuning (number of estimators, maximum depth, learning rate,
minimum child weight, and training-weight configuration) supporting model
selection on validation performance, with PR-AUC given particular attention
because of the class imbalance.

Evaluation uses accuracy, precision, recall, F1-score, ROC-AUC, and PR-AUC,
together with a confusion matrix for error analysis. SHAP is used to generate
both local explanations (for individual predictions) and global feature-importance
analysis.

\section{Provision and confirmation of the results}

The component report states that the current results are provisional:
93.00\% accuracy, 86.00\% precision, 89.00\% recall, 87.47\% F1-score, 97.20\%
ROC-AUC, and 90.30\% PR-AUC. These figures were placed in the draft before the
confirmed, verified evaluation of the final model configuration was earned.
They are reproduced in \cref{ch:results} with the same caveat: final
conclusions require the confirmed evaluation results.

\section{Component artefacts}

\begin{description}
  \item[Scanner-output parsers] -- for Nmap, Nikto, and Gobuster execution
  results.
  \item[Intelligence enricher] -- applies NVD, KEV, and EPSS enrichment.
  \item[Trained classifier] -- the XGBoost model over the five features.
  \item[SHAP explainer] -- local and global explanations.
  \item[Neo4j graph] -- persistence of sessions, hosts, services, and
  vulnerabilities.
  \item[FastAPI interface] -- the \texttt{/analyze} endpoint that consumes
  structured execution results and returns the analysis output.
\end{description}